\documentclass[11pt]{article}
\usepackage[margin=1in]{geometry}
\usepackage{amsmath,amssymb,amsthm}
\usepackage{graphicx}
\usepackage[hidelinks]{hyperref}
\usepackage{microtype}

\newtheorem{theorem}{Theorem}
\newtheorem{lemma}{Lemma}
\newcommand{\R}{\mathbb R}
\newcommand{\linf}{\ell_\infty}

\title{Two Closed Subspace-Cones Refute Both\\
Lipschitz Klee--And\^o Extensions}
\author{Candidate solution packet for arXiv:1606.00249}
\date{9 August 2026}

\begin{document}
\maketitle

\begin{center}
\fbox{\parbox{0.91\linewidth}{\textbf{Status.} Candidate counterexample,
likely valid: a full negative answer to Problem 1.5 of
arXiv:1606.00249, subject to human review.}}
\end{center}

\section{Source question}

Messerschmidt's paper \cite{source} proves continuous positively homogeneous
selection theorems for arbitrary collections of closed cones.  On page 3 it
asks:

\begin{quote}
\textbf{Problem 1.5.} Do Corollaries 5.5 and 5.6 remain true when, in
their statements, the word ``continuous'' is replaced by the word
``Lipschitz''?
\end{quote}

Here Corollary 5.5 is the coadditivity theorem and Corollary 5.6 is the
conormality theorem.  The source defines a cone as a nonempty set closed under
addition and multiplication by nonnegative scalars; properness is a separate
condition.  Thus every linear subspace is a cone.

\begin{figure}[ht]
  \centering
  \includegraphics[width=\linewidth]{figures/open_problem_crop.png}
  \caption{Problem 1.5 in \cite[p.~3]{source}.}
\end{figure}

\section{Definitions and the quotient obstruction}

Let $q:E\to Y$ be a bounded linear surjection of real Banach spaces.  A right
inverse is a map $s:Y\to E$ satisfying $q\circ s=\mathrm{id}_Y$; it need not be
linear.  We use the classical fact that
\[
  q:\linf\longrightarrow \linf/c_0
\]
has no uniformly continuous right inverse, and hence no Lipschitz right
inverse.  This is Kalton's example, recorded as Proposition 2.5 (and again in
Corollary 5.9) of \cite{quotient}.

For two cones $C_1,C_2\subseteq X$, the conormality selection requested by the
Lipschitz version of Corollary 5.6 consists of Lipschitz maps
$\delta_i:X\to C_i$ with
\[
  x=\delta_1(x)+\delta_2(x).
\]
The coadditivity selection requested by the Lipschitz version of Corollary 5.5
is a Lipschitz map $v:X^2\to X$ such that
\[
  v(\xi_1,\xi_2)\in(\xi_1+C_1)\cap(\xi_2+C_2).
\]
For a two-point index set the source's spaces $c(\Omega,X)$ and
$\ell_\infty(\Omega,X)$ both equal $X^2$, with the maximum norm.

\section{Proof intuition}

The graph of a quotient map turns lifting into cone geometry.  Decomposing
$(0,y)$ into a horizontal vector and a graph vector necessarily chooses an
$e\in E$ with $q(e)=y$.  Similarly, an intersection point between a horizontal
translate and a graph translate chooses such an $e$, up to sign.  Therefore a
Lipschitz selector in either formulation would give a Lipschitz right inverse
of $q$.  Choosing a quotient known not to have such a right inverse produces
one pair of closed cones that defeats both proposed extensions.

\section{Counterexample theorem}

\begin{theorem}\label{thm:main}
There exist a Banach space $X$ and two closed cones $C_1,C_2\subseteq X$ for
which all of the following hold.
\begin{enumerate}
  \item Every $x\in X$ can be written $x=c_1+c_2$ with $c_i\in C_i$.
  \item For every $(\xi_1,\xi_2)\in X^2$, the intersection
  $(\xi_1+C_1)\cap(\xi_2+C_2)$ is nonempty.
  \item There are no Lipschitz maps $\delta_i:X\to C_i$ satisfying
  $x=\delta_1(x)+\delta_2(x)$ for every $x\in X$.
  \item There is no Lipschitz map $v:X^2\to X$ satisfying
  $v(\xi_1,\xi_2)\in(\xi_1+C_1)\cap(\xi_2+C_2)$ for every
  $(\xi_1,\xi_2)\in X^2$.
\end{enumerate}
Consequently, both clauses of Problem 1.5 in \cite{source} have negative
answers.  The failure persists even if positive homogeneity and the
quantitative norm bounds are omitted from the desired conclusions.
\end{theorem}

\begin{proof}
Let $q:E\to Y$ be any bounded linear surjection having no Lipschitz right
inverse.  Equip
\[
  X:=E\oplus_\infty Y
\]
with its maximum norm, and define
\[
  M:=E\oplus\{0\},\qquad
  N:=\operatorname{graph}(q)=\{(e,qe):e\in E\}.
\]
Both $M$ and $N$ are closed linear subspaces of $X$, hence closed cones in the
sense of \cite{source}.  Set $C_1=M$ and $C_2=N$.

First, $M+N=X$.  Indeed, for $(a,y)\in X$, choose $e\in E$ with $qe=y$.
Then
\[
  (a,y)=(a-e,0)+(e,qe)\in M+N.
\]
This proves assertion 1.  It also proves assertion 2: given
$\xi_1,\xi_2\in X$, write $\xi_2-\xi_1=m+n$ with $m\in M$ and $n\in N$.
Then
\[
  \xi_1+m=\xi_2-n\in(\xi_1+M)\cap(\xi_2+N).
\]

Suppose assertion 3 fails.  Thus Lipschitz maps $\delta_M:X\to M$ and
$\delta_N:X\to N$ satisfy
\[
  z=\delta_M(z)+\delta_N(z)\qquad(z\in X).
\]
For $y\in Y$, write
\[
  \delta_N(0,y)=(s(y),y).
\]
The second coordinate is $y$ because the second coordinate of
$\delta_M(0,y)$ is zero.  Since $\delta_N(0,y)\in N$, we have $q(s(y))=y$.
Moreover $s$ is Lipschitz, being the composition of the bounded embedding
$y\mapsto(0,y)$, the Lipschitz map $\delta_N$, and the bounded first-coordinate
projection.  Thus $s$ is a Lipschitz right inverse of $q$, a contradiction.

Suppose assertion 4 fails, and let $v$ be the asserted Lipschitz selector.  For
$y\in Y$, the first membership condition gives
\[
  v\bigl(0,(0,y)\bigr)\in M,
\]
so write $v(0,(0,y))=(a(y),0)$.  The second membership condition gives
\[
  (a(y),-y)=v\bigl(0,(0,y)\bigr)-(0,y)\in N.
\]
Hence $q(a(y))=-y$.  The map $s(y):=-a(y)$ is Lipschitz and satisfies
$q(s(y))=y$, again contradicting the choice of $q$.

Finally take $E=\linf$, $Y=\linf/c_0$, and let $q$ be the quotient map.  By
\cite[Proposition~2.5]{quotient}, this $q$ has no uniformly continuous right
inverse, so in particular no Lipschitz right inverse.  The preceding
construction proves all four assertions.  Neither contradiction used positive
homogeneity or a norm estimate.
\end{proof}

\section{Verification notes and scope}

The proof is structural and uses no computation.  The points most useful to
check are:
\begin{itemize}
  \item $N$ is closed because it is the graph of a bounded linear map;
  \item subspaces qualify as cones under the source definition;
  \item restricting either hypothetical selector to the displayed copy of
  $Y$ really produces a right inverse of $q$;
  \item a two-element index set simultaneously covers both allowed choices of
  $Z$ in Corollary 5.5.
\end{itemize}

This theorem settles Problem 1.5 as stated, because that problem concerns the
source's arbitrary-collection results.  It does not settle the separate,
stronger Conjecture 1.6, which requests one closed generating cone and
Lipschitz positive and negative parts for that single cone.

\section{Bounded novelty check}

On 9 August 2026 we searched the run indexes, the source text, exact phrases
from Problem 1.5, and close variants involving Klee--And\^o selections,
Lipschitz decomposition, and quotient right inverses.  We also checked the
2019 survey on the single-cone Lipschitz decomposition problem and
arXiv:1909.10417.  We found the known quotient obstruction used above, but no
source stating the two-subspace graph reduction or claiming a negative answer
to both Corollaries 5.5 and 5.6.  This is a bounded novelty check, not an
exhaustive bibliographic certification.

\section{Human-review recommendation}

Recommended verdict: likely valid candidate counterexample.  Review should
focus on the two restrictions that extract a right inverse and on the scope
distinction between Problem 1.5 and Conjecture 1.6.

\begin{thebibliography}{9}
\bibitem{source}
M.~Messerschmidt,
\emph{Strong Klee--And\^o Theorems through an Open Mapping Theorem for
cone-valued multi-functions}, arXiv:1606.00249v2, 2018.

\bibitem{quotient}
M.~Messerschmidt,
\emph{A family of quotient maps of $\ell^\infty$ that do not admit uniformly
continuous right inverses}, arXiv:1909.10417v1, 2019.
\end{thebibliography}

\end{document}
